\task{Dreiecksschema und Fragen}
{
	Betrachten wir nun die Interpolation mit einem Dreiecksschema und gucken uns etwas Verständnis an.
	}{
		\itemf Gegeben seien die Stützstellen einer uns unbekannten Funktion $f(x)$: 
		\begin{align*}
			f(0) = 3; \quad f(1) = 2; \quad f(8) = 2
		\end{align*}
		Berechnen Sie die Interpolation mithilfe des Newton-Verfahrens und stellen Sie das Dreiecksschema auf.
		\itemf Bei gleichbleibenden Stützstellen, wie unterscheidet sich (ganz allgemein) das Ergebnis einer Interpolation mit den Lagrange-Basisfunktionen gegenüber dem eines Newton-Verfahrens?
		\itemf Gilt bei Polynominterpolation immer "Je mehr Stützstellen, desto höher ist die Genauigkeit der Interpolation"?
		\itemf Welcher Vorteil bietet eine stückweise Interpolation gegenüber einer nicht-stückweisen?
		% aitken-neville
		% dreiecksschema
		% 
	}{
		\itemf Stellen wir zuerst das Dreiecksschema auf:
		\begin{table}[H]
			\centering
			\begin{tabular}{cc|cccc}
				\toprule
				$x_i$  & $i \setminus k$ & 0 & 1 & 2 \\ \midrule
				$0$ & $0$ & $3$ & $c_{0,1}$ & $c_{0,2}$ \\
				$1$ & $1$ & $2$ & $c_{1,1}$ & \\
				$8$ & $2$ & $2$ & &
			\end{tabular}
		\end{table}%
		\begin{align*} 
			c_{0,1} &= \frac{2 - 3}{1 - 0} = -1 \\
			c_{1,1} &= \frac{2 - 2}{8 - 1} = 0 \\
			c_{0,2} &= \frac{0 + 1}{8 - 0} = \frac{1}{8}
		\end{align*}
		\begin{table}[H]
			\centering
			\begin{tabular}{cc|cccc}
				\toprule
				$x_i$  & $i \setminus k$ & 0 & 1 & 2 \\ \midrule
				$0$ & $0$ & $3$ & $-1$ & $\nicefrac{1}{8}$ \\
				$1$ & $1$ & $2$ & $0$ & \\
				$8$ & $2$ & $2$ & & \\
			\end{tabular}
		\end{table}
		Jetzt noch die Interpolationsfunktion aufstellen:
		\begin{align*}
			p(x) &= 3 + (-1) \cdot (x-x_0) + \frac{1}{8} \cdot (x-x_0)(x-x_1) \\
				&= 3 + (-1) \cdot (x-0) + \frac{1}{8} \cdot (x-0)(x-1) \\
				&= 3 + (-x) + \frac{1}{8} \cdot (x^2-x) \\
				&= \frac{1}{8}x^2 - \frac{9}{8}x + 3
		\end{align*}
		\itemf Es unterscheidet sich gar nicht, denn Interpolationspolynome sind stets eindeutig!
		\itemf Theoretisch erstmal schon, aber bei vielen Stützstellen kommt der Runge-Effekt zum Vorschein und macht Ergebnisse deutlich ungenauer.
		\itemf Stückweise Interpolation wirkt dem Runge-Effekt entgegen, da man seine Anzahl Stützstellen in Stücke unterteilt (und damit jedes Stück mit weniger Stützstellen arbeitet und somit der Runge-Effekt schlechter auftreten kann).
}